\section{Hypothesis}

This research hypothesises that representing detected joints as nodes in a
depth-augmented graph and processing this graph with a GAT will improve both
the accuracy and the computational efficiency of multi-person pose tracking
compared with approaches that rely solely on 2-D image coordinates.

Attaching a monocular depth value to every joint lifts the skeleton into
2.5-D space, giving the GAT an additional geometric cue that helps distinguish
overlapping limbs and nearby individuals.  Message passing allows each node to
aggregate information from its neighbours, while the attention mechanism
assigns larger weights to edges whose combined spatial and depth attributes
suggest that the connected joints belong to the same person.  

Because the graph encodes only compact descriptors -- per-joint image
coordinates, depth, confidence, a joint-type embedding, and a small set of
edge attributes such as inter-joint distance, depth difference, and local
orientation -- its feature space is far smaller than the dense heat-maps or
pixel-level tensors used by many convolutional pipelines.  This compact
representation lowers memory usage, shortens training time, and enables
real-time inference on modest hardware.  At the same time, the depth channel
and geometric edge cues supply context that pure 2-D coordinates lack,
allowing the model to maintain or even improve joint-to-person assignment
accuracy without resorting to deeper or wider backbones.